% =============================================================================
% 4 Konzeptentwicklung
% =============================================================================
\section{Konzeptentwicklung}
\label{chap:konzeptentwicklung}

Die Analyse des bestehenden Erwärmungsprüfstands hat gezeigt, dass die
zentrale Regelung zwar die unmittelbar hinter dem Festtransformator
gemessenen Phasenströme nachführt, die Verteilung dieser Ströme auf die
einzelnen Abgänge des MCC-Feldes jedoch nicht festlegt. Ändern sich die
Widerstände der angeschlossenen Lastzweige infolge ihrer Erwärmung,
weichen daher die einzelnen Modulströme von ihren Sollwerten ab. Diese
Abweichungen müssen im bestehenden Prüfablauf durch manuelles
Verstellen der V2A-Widerstände ausgeglichen werden.

Ausgehend von den in
Abschnitt~\ref{sec:anforderungen_optimierung}
abgeleiteten Anforderungen wird im Folgenden ein
transformatorgestütztes Stellkonzept entwickelt. Ziel ist die
individuelle und automatisierte Nachführung jedes Phasenstroms der
zwölf dreiphasigen MCC-Abgänge. Die vorhandene Prüfeinspeisung aus
Säulenstelltransformator, Festtransformator, Leistungsschalterfeld,
Anbindefeld und MCC-Feld soll dabei grundsätzlich erhalten bleiben.

Als elektrisches Stellglied wird für jeden Phasenstrom ein zusätzlicher
Einphasentransformator eingesetzt, der im Folgenden als Trafo~2
bezeichnet wird. An dessen Sekundärwicklung wird ein motorisch
verstellbarer Widerstand \(R_{\mathrm{Stell}}\) angeschlossen. Durch
die Impedanztransformation erscheint dieser Widerstand im
niederspannungsseitigen Hochstromkreis mit einem wesentlich kleineren
Wert. Die eigentliche Widerstandsverstellung kann dadurch auf der
stromärmeren Sekundärseite erfolgen.

Die Entwicklung umfasst in diesem Kapitel das Wirkprinzip, die
Anordnung der Stellkanäle, die vorgesehene Mess- und Regelungsstruktur
sowie eine überschlägige Vorbemessung. Die detaillierte Abbildung der
elektrischen und thermischen Zusammenhänge in MATLAB/Simulink ist
Gegenstand von Kapitel~\ref{chap:modellbildung}.


% -----------------------------------------------------------------------------
% 4.1 Entwicklungsziel und Systemarchitektur
% -----------------------------------------------------------------------------
\subsection{Entwicklungsziel und Systemarchitektur}
\label{sec:entwicklungsziel_systemarchitektur}

Das Stellkonzept muss die Phasenströme der sechs
\SI{40}{\ampere}-Abgänge, der drei \SI{70}{\ampere}-Abgänge sowie der
Abgänge mit \SI{250}{\ampere}, \SI{400}{\ampere} und
\SI{630}{\ampere} voneinander unabhängig beeinflussen können. Da
sämtliche Abgänge dreiphasig betrieben werden, ist für jeden Abgang
und jeden Außenleiter ein eigener Stellkanal erforderlich.

Die Anzahl der Stellkanäle beträgt somit

\begin{equation}
    n_{\mathrm{Stell}}
    =
    12 \cdot 3
    =
    36.
    \label{eq:anzahl_stellkanaele}
\end{equation}

Das Konzept umfasst daher insgesamt 36 Transformatoren des Typs
Trafo~2 und 36 zugehörige motorisch verstellbare Widerstände. Die
Aufteilung der Stellkanäle ist in
Tabelle~\ref{tab:aufteilung_stellkanaele} zusammengefasst.

\begin{table}[H]
    \centering
    \caption{Aufteilung der Stellkanäle auf die MCC-Abgänge}
    \label{tab:aufteilung_stellkanaele}
    \begin{tabular}{cccc}
        \toprule
        Prüfstrom je Abgang &
        Anzahl Abgänge &
        Stellkanäle je Abgang &
        Anzahl Stellkanäle \\
        \midrule
        \SI{40}{\ampere}  & 6 & 3 & 18 \\
        \SI{70}{\ampere}  & 3 & 3 & 9  \\
        \SI{250}{\ampere} & 1 & 3 & 3  \\
        \SI{400}{\ampere} & 1 & 3 & 3  \\
        \SI{630}{\ampere} & 1 & 3 & 3  \\
        \midrule
        Gesamt            & 12 & -- & 36 \\
        \bottomrule
    \end{tabular}
\end{table}

Für die weitere Beschreibung bezeichnet der Index \(k\) den
betrachteten MCC-Abgang mit

\begin{equation}
    k \in \left\{1,\ldots,12\right\}.
    \label{eq:definition_abgangsindex}
\end{equation}

Der Index \(p\) kennzeichnet den zugehörigen Außenleiter:

\begin{equation}
    p \in
    \left\{
        \mathrm{L1},
        \mathrm{L2},
        \mathrm{L3}
    \right\}.
    \label{eq:definition_phasenindex}
\end{equation}

Für die Abgänge mit \SI{250}{\ampere}, \SI{400}{\ampere} und
\SI{630}{\ampere} bleibt der vorhandene V2A-Widerstand als
Grundlastpfad erhalten. Trafo~2 bildet dort einen parallel
angeordneten Zusatzpfad, über den der zur Einstellung und Nachführung
benötigte Stromanteil geführt wird.

Bei den Abgängen mit \SI{40}{\ampere} und \SI{70}{\ampere} entfällt
der V2A-Grundlastpfad. Der vollständige Phasenstrom wird hier durch
die Primärwicklung von Trafo~2 geführt und über den sekundärseitigen
Stellwiderstand eingestellt.

Die Stromaufteilung beziehungsweise Strombeeinflussung erfolgt jeweils
hinter dem zu prüfenden MCC-Modul. Der vollständige Modulstrom
durchfließt somit zunächst die in der Abgangseinheit enthaltenen
Schaltgeräte, Kontakte und Strombahnen. Erst anschließend wird der
Strom über den jeweiligen Lastpfad zum gemeinsamen Sternpunkt
geführt. Die für die Erwärmungsprüfung notwendige Belastung des
MCC-Moduls bleibt dadurch erhalten.

\begin{center}
    \fbox{
        \parbox{0.90\textwidth}{
            \centering
            \textcolor{red}{
                Abbildung ergänzen: Gesamtstruktur des
                transformatorgestützten Stellkonzepts mit
                zwölf dreiphasigen MCC-Abgängen,
                36 Stellkanälen und gemeinsamem Sternpunkt.
            }
        }
    }
\end{center}

% Vorgesehener Dateiname:
% 03_Ressourcen/drawio/gesamtstruktur_stellkonzept.pdf


% -----------------------------------------------------------------------------
% 4.2 Wirkprinzip des transformatorgestützten Stellkanals
% -----------------------------------------------------------------------------
\subsection{Wirkprinzip des transformatorgestützten Stellkanals}
\label{sec:wirkprinzip_trafo2}

Ein Stellkanal besteht aus der Primärwicklung von Trafo~2, dessen
Sekundärwicklung, dem motorisch verstellbaren Widerstand
\(R_{\mathrm{Stell},k,p}\), der zugehörigen Strommessung und einem
Regelkanal. Trafo~2 stellt dabei keine zusätzliche Energiequelle dar.
Die im Stellwiderstand umgesetzte Leistung wird weiterhin durch die
bestehende Hochstromeinspeisung bereitgestellt.

Die Aufgabe des Transformators besteht darin, die auf seiner
Sekundärseite angeschlossene Impedanz auf einen für den
niederspannungsseitigen Hochstromkreis geeigneten Wert zu
transformieren. Für den idealisierten Transformator gilt das
Übersetzungsverhältnis

\begin{equation}
    u_{\mathrm{T2}}
    =
    \frac{N_1}{N_2}
    =
    \frac{U_{\mathrm{T2},1}}{U_{\mathrm{T2},2}}
    =
    \frac{I_{\mathrm{T2},2}}{I_{\mathrm{T2},1}}.
    \label{eq:uebersetzungsverhaeltnis_trafo2}
\end{equation}

Dabei bezeichnen \(N_1\) und \(N_2\) die Windungszahlen,
\(U_{\mathrm{T2},1}\) und \(U_{\mathrm{T2},2}\) die Effektivwerte der
Spannungen sowie \(I_{\mathrm{T2},1}\) und
\(I_{\mathrm{T2},2}\) die Effektivwerte der Ströme auf der Primär-
und Sekundärseite. Die grundlegenden Transformatorbeziehungen und die
Übertragung von Impedanzen werden in
\textcite[Abschn.~2.1.2--2.1.5 und 2.3.3,
S.~109--114 und 136--141]{spring2009maschinen} hergeleitet.

Trafo~2 soll die niedrige Spannung des Hochstromkreises auf ein
höheres sekundärseitiges Spannungsniveau transformieren. Damit gilt

\begin{equation}
    N_2 > N_1
    \qquad\Longrightarrow\qquad
    u_{\mathrm{T2}} < 1.
    \label{eq:trafo2_spannungserhoehung}
\end{equation}

Der sekundärseitige Stellwiderstand erscheint auf der Primärseite mit

\begin{equation}
    R_{\mathrm{Stell},k,p}'
    =
    u_{\mathrm{T2}}^{\,2}
    R_{\mathrm{Stell},k,p}.
    \label{eq:stellwiderstand_transformation}
\end{equation}

Da \(u_{\mathrm{T2}}<1\) ist, wird der sekundärseitige Widerstand auf
einen deutlich kleineren primärseitig wirksamen Widerstand
transformiert. Dadurch kann ein Stellwiderstand im Ohmbereich einen
Hochstromkreis beeinflussen, dessen wirksame Widerstände im
Milliohmbereich liegen.

Eine Verringerung von \(R_{\mathrm{Stell},k,p}\) reduziert dessen
primärseitig wirksamen Wert. Bei gleichbleibender Speisespannung
steigt dadurch der Primärstrom von Trafo~2:

\begin{equation}
    R_{\mathrm{Stell},k,p} \downarrow
    \quad\Longrightarrow\quad
    R_{\mathrm{Stell},k,p}' \downarrow
    \quad\Longrightarrow\quad
    I_{\mathrm{T2},1,k,p} \uparrow.
    \label{eq:stellrichtung_trafo2}
\end{equation}

Für eine Vergrößerung des Stellwiderstands ergibt sich die
entgegengesetzte Wirkung. Der Stellwiderstand kann somit als
motorisch betätigtes Stellglied der Stromregelung verwendet werden.

Die erreichbare Stromänderung wird im realen Aufbau zusätzlich durch
die Wicklungswiderstände und Streureaktanzen von Trafo~2 begrenzt.
Insbesondere die Kurzschlussimpedanz beziehungsweise die relative
Kurzschlussspannung bestimmt den maximal erreichbaren Strom bei
kleinem sekundärseitigem Stellwiderstand. Die physikalische Bedeutung
dieser Größen wird bei
\textcite[Abschn.~2.3.4--2.4.2,
S.~142--154]{spring2009maschinen} beschrieben. Ihre detaillierte
Berücksichtigung erfolgt im elektrischen Simulationsmodell in
Kapitel~\ref{chap:modellbildung}.


% -----------------------------------------------------------------------------
% 4.3 Ausführung für 250-A-, 400-A- und 630-A-Abgänge
% -----------------------------------------------------------------------------
\subsection{Ausführung für die Abgänge mit hohen Prüfströmen}
\label{sec:ausfuehrung_hohe_modulstroeme}

Bei den Abgängen mit Prüfströmen von \SI{250}{\ampere},
\SI{400}{\ampere} und \SI{630}{\ampere} bleiben die vorhandenen
V2A-Widerstände als Grundlastpfade erhalten. Für jeden Außenleiter
wird die Primärwicklung eines eigenen Trafo~2 parallel zum
zugehörigen V2A-Widerstand geschaltet.

Der gesamte Phasenstrom durchfließt zunächst das MCC-Modul und teilt
sich anschließend auf den V2A-Grundlastpfad und den
Transformatorzweig auf. Beide Teilströme werden am gemeinsamen
Sternpunkt wieder zusammengeführt.

Da der Transformatorzweig neben einem ohmschen auch einen induktiven
Anteil besitzt, ist die Stromaufteilung grundsätzlich als
Zeigerbeziehung zu beschreiben:

\begin{equation}
    \underline{I}_{\mathrm{MCC},k,p}
    =
    \underline{I}_{\mathrm{V2A},k,p}
    +
    \underline{I}_{\mathrm{T2},1,k,p}.
    \label{eq:stromaufteilung_hochstromabgang}
\end{equation}

Der V2A-Widerstand übernimmt den überwiegenden Anteil des
Modulstroms. Trafo~2 muss deshalb nicht für den vollständigen
Prüfstrom des Abgangs, sondern für den erforderlichen Grund-,
Stell- und Nachführstrom ausgelegt werden. Hierdurch lässt sich die
Bemessungsleistung des zusätzlichen Transformators gegenüber einer
vollständigen Stromführung reduzieren.

Vor Beginn der Prüfung wird der V2A-Widerstand so eingestellt, dass
sein Teilstrom unterhalb des geforderten Modulstroms liegt. Die
verbleibende Differenz wird bereits im kalten Ausgangszustand durch
Trafo~2 bereitgestellt:

\begin{equation}
    I_{\mathrm{T2},1,k,p,\mathrm{kalt}}
    =
    I_{\mathrm{soll},k,p}
    -
    I_{\mathrm{V2A},k,p,\mathrm{kalt}}.
    \label{eq:trafo2_grundstrom_hochstromabgang}
\end{equation}

Diese Einstellung ist erforderlich, da der Transformatorzweig einen
zusätzlichen Strombeitrag bereitstellen, jedoch keinen Strom aus dem
V2A-Pfad abziehen kann. Der V2A-Grundlastpfad darf den Sollstrom daher
auch bei kaltem Widerstand nicht überschreiten.

Mit zunehmender Temperatur steigt der Widerstand der V2A-Schienen.
Bei näherungsweise konstanter Phasenspannung sinkt dadurch der
V2A-Teilstrom. Der Stellwiderstand wird daraufhin verkleinert, sodass
der Primärstrom von Trafo~2 zunimmt und der thermisch bedingte
Stromrückgang ausgeglichen wird:

\begin{equation}
    I_{\mathrm{T2},1,k,p,\mathrm{warm}}
    =
    I_{\mathrm{soll},k,p}
    -
    I_{\mathrm{V2A},k,p,\mathrm{warm}}.
    \label{eq:trafo2_nachfuehrstrom_hochstromabgang}
\end{equation}

Für die drei Abgänge mit hohen Prüfströmen werden jeweils drei
phasenselektive Stellkanäle benötigt. Diese Ausführungsvariante
umfasst somit neun Transformatoren und neun motorisch verstellbare
Widerstände.


% -----------------------------------------------------------------------------
% 4.4 Ausführung für 40-A- und 70-A-Abgänge
% -----------------------------------------------------------------------------
\subsection{Ausführung für die Abgänge mit kleinen Prüfströmen}
\label{sec:ausfuehrung_kleine_stroeme}

Bei den sechs Abgängen mit jeweils \SI{40}{\ampere} und den drei
Abgängen mit jeweils \SI{70}{\ampere} wird kein paralleler
V2A-Grundlastpfad eingesetzt. Der vollständige Phasenstrom fließt
durch die Primärwicklung des jeweils zugeordneten Trafo~2:

\begin{equation}
    \underline{I}_{\mathrm{MCC},k,p}
    =
    \underline{I}_{\mathrm{T2},1,k,p}.
    \label{eq:strom_kleinabgang_trafo2}
\end{equation}

Der primärseitige Transformatorstrom wird durch den auf die
Primärseite transformierten Stellwiderstand und die innere Impedanz
von Trafo~2 bestimmt. Der Transformatorzweig dient bei diesen
Abgängen daher nicht nur zur Korrektur einer Stromabweichung, sondern
bildet den vollständigen Lastpfad des jeweiligen Außenleiters.

Der Stellbereich muss sowohl die erstmalige Einstellung des
Sollstroms als auch dessen Nachführung während der Erwärmungsprüfung
ermöglichen. Neben der Eigenerwärmung von Trafo~2 und
\(R_{\mathrm{Stell}}\) sind dabei auch die temperaturbedingten
Änderungen der Modul-, Kontakt- und Verbindungswiderstände zu
berücksichtigen.

Für die neun kleinen MCC-Abgänge werden insgesamt 27 Stellkanäle
benötigt. Eine identische Dimensionierung der
\SI{40}{\ampere}- und \SI{70}{\ampere}-Kanäle ist nicht zwingend
erforderlich. Das Wirkprinzip und die regelungstechnische Ansteuerung
bleiben jedoch gleich.

\begin{center}
    \fbox{
        \parbox{0.90\textwidth}{
            \centering
            \textcolor{red}{
                Abbildung ergänzen: Gegenüberstellung eines
                Hochstromabgangs mit parallelem V2A- und
                Trafo-2-Pfad sowie eines 40-A-/70-A-Abgangs mit
                vollständiger Stromführung durch Trafo~2.
            }
        }
    }
\end{center}

% Vorgesehener Dateiname:
% 03_Ressourcen/drawio/stellkanaele_trafo2.pdf


% -----------------------------------------------------------------------------
% 4.5 Mess- und Regelungskonzept
% -----------------------------------------------------------------------------
\subsection{Mess- und Regelungskonzept}
\label{sec:mess_regelungskonzept}

Im bestehenden Prüfstand werden die hinter dem Festtransformator
gemessenen Phasenströme als Regelgrößen des Säulenstelltransformators
verwendet. Diese übergeordnete Stromregelung ist nicht mit einer
gleichzeitigen individuellen Regelung der 36 Modulströme vereinbar.

Erhöht ein nachgelagerter Stellkanal den Strom eines einzelnen
MCC-Abgangs, steigt gleichzeitig der von der gemeinsamen Einspeisung
bereitgestellte Phasenstrom. Die bestehende Gesamtstromregelung würde
darauf mit einer Verringerung der Ausgangsspannung reagieren. Diese
Spannungsänderung würde wiederum sämtliche parallel angeschlossenen
MCC-Abgänge beeinflussen. Die übergeordnete Gesamtstromregelung und
die nachgelagerten Modulstromregelungen würden folglich gegeneinander
arbeiten.

Aus diesem Grund wird die vorgelagerte Stromregelung durch eine
phasenselektive Spannungsregelung ersetzt. Der
Säulenstelltransformator regelt zukünftig die drei Phasenspannungen
an der gemeinsamen Niederspannungssammelschiene. Die 36
nachgelagerten Regelkreise stellen anschließend die einzelnen
Modulströme über die zugehörigen Stellwiderstände ein.

Für die vorgelagerte Spannungsregelung ergibt sich je Außenleiter die
Regelabweichung

\begin{equation}
    e_{U,p}
    =
    U_{\mathrm{NS,soll},p}
    -
    U_{\mathrm{NS,ist},p}.
    \label{eq:regelabweichung_phasenspannung}
\end{equation}

Als Stellglieder dienen weiterhin die motorisch verfahrbaren
Stromabnehmer des Säulenstelltransformators. Die Regelung der drei
Phasenspannungen ist notwendig, da die Stellungen der drei
Phasenantriebe während des Betriebs voneinander abweichen können.

Der bisher als Regelgröße verwendete übergeordnete Phasenstrom wird
weiterhin erfasst, jedoch nicht mehr durch den
Säulenstelltransformator auf einen festen Wert geregelt. Er ergibt
sich aus dem Strom des Leistungsschalterpfads und der Summe der
parallel angeschlossenen MCC-Abgänge:

\begin{equation}
    \underline{I}_{\mathrm{ges},p}
    =
    \underline{I}_{\mathrm{KB},p}
    +
    \sum_{k=1}^{12}
    \underline{I}_{\mathrm{MCC},k,p}.
    \label{eq:gesamtstrom_neues_regelkonzept}
\end{equation}

Die Messung des Gesamtstroms bleibt für die Überwachung der
Belastungsgrenzen des Festtransformators, der Kupferschienen, des
Anbindefeldes und des Leistungsschalterpfads erforderlich.

Für jeden MCC-Abgang und Außenleiter wird ein eigener Stromregelkreis
vorgesehen. Der Stromsensor muss jeweils den vollständigen Strom durch
das MCC-Modul erfassen. Bei den hohen Prüfströmen ist die Messstelle
daher vor der Aufteilung auf den V2A- und Transformatorpfad
anzuordnen.

Die Regelabweichung eines Stromkanals lautet

\begin{equation}
    e_{I,k,p}
    =
    I_{\mathrm{soll},k,p}
    -
    I_{\mathrm{ist},k,p}.
    \label{eq:regelabweichung_modulstrom}
\end{equation}

Ist der gemessene Modulstrom kleiner als sein Sollwert, wird
\(R_{\mathrm{Stell},k,p}\) verkleinert. Der Primärstrom von Trafo~2
steigt. Überschreitet der Modulstrom seinen Sollwert, wird der
Stellwiderstand vergrößert.

Die Trennung in eine gemeinsame Spannungsregelung und dezentrale
Stromregelkreise folgt dem regelungstechnischen Grundsatz, eine
komplexe Strecke in eindeutig zugeordnete Teilregelkreise zu
unterteilen. Die grundlegende Beschreibung geschlossener
Regelkreise anhand von Führungs-, Rückführungs- und Störgrößen ist bei
\textcite[Abschn.~3.4, S.~96--102]
{schroeder2021elektrische_antriebe} dargestellt. Vorteile der
Unterteilung komplexer Regelstrecken in einfachere Teilregelkreise
werden von
\textcite[Abschn.~4.3, S.~153--155]
{schroeder2021elektrische_antriebe} erläutert.

Bei der entwickelten Struktur handelt es sich jedoch nicht um eine
klassische Kaskadenregelung, da die Sollwerte der Modulstromregler
nicht durch den Spannungsregler erzeugt werden. Die Spannungsregelung
stellt vielmehr eine möglichst konstante gemeinsame elektrische
Randbedingung für die nachgelagerten Stromregelkreise bereit.

Vollständig unabhängig sind die 36 Stromregelkreise dennoch nicht.
Alle Abgänge werden über dieselben Niederspannungssammelschienen
gespeist. Eine Änderung eines Stellwiderstands verändert die Belastung
der gemeinsamen Quelle und kann infolge der endlichen
Quellen- und Leitungsimpedanzen andere Abgänge beeinflussen.
Zusätzlich bestehen Kopplungen zwischen den drei Außenleitern der
Transformatoren. Der Prüfstand ist daher auch mit der neuen
Regelungsstruktur als gekoppeltes Mehrgrößensystem zu behandeln.

Durch die vorgelagerte Spannungsregelung wird diese Kopplung nicht
vollständig beseitigt. Sie verhindert jedoch, dass die zentrale
Regelung jede beabsichtigte Änderung eines Modulstroms unmittelbar
durch eine gegenläufige Veränderung der gemeinsamen Speisespannung
kompensiert.

Beim Einschalten werden die Stellwiderstände zunächst in eine
strombegrenzende Ausgangsposition gefahren. Danach werden die
Phasenspannungen über den Säulenstelltransformator aufgebaut. Erst
nach Erreichen eines stabilen Spannungsarbeitspunkts werden die
nachgelagerten Stromregelkreise freigegeben.

Die Soll- und Istwerte der Phasenspannungen, die übergeordneten
Phasenströme, die 36 Modulströme und die Positionen der
Stellwiderstände sollen durch die SPS erfasst und in WinCC dargestellt
und aufgezeichnet werden. Zusätzlich sind die elektrischen und
thermischen Belastungsgrenzen der Stelltransformatoren und
Stellwiderstände zu überwachen.

\begin{center}
    \fbox{
        \parbox{0.90\textwidth}{
            \centering
            \textcolor{red}{
                Abbildung ergänzen: Regelungsarchitektur mit drei
                vorgelagerten Spannungsregelkreisen, gemeinsamer
                Niederspannungssammelschiene und 36 dezentralen
                Modulstromregelkreisen.
            }
        }
    }
\end{center}

% Vorgesehener Dateiname:
% 03_Ressourcen/drawio/regelstruktur_stellkonzept.pdf


% -----------------------------------------------------------------------------
% 4.6 Vorbemessung und erforderliche Stellreserven
% -----------------------------------------------------------------------------
\subsection{Vorbemessung und erforderliche Stellreserven}
\label{sec:dimensionierung_stellreserve}

Die prinzipielle Funktionsfähigkeit des Konzepts setzt voraus, dass
der Sollstrom jedes Stellkanals innerhalb seines technisch
erreichbaren Strombereichs liegt. Für jeden Abgang und jeden
Außenleiter muss daher gelten:

\begin{equation}
    I_{\mathrm{MCC},k,p,\min}
    <
    I_{\mathrm{soll},k,p}
    <
    I_{\mathrm{MCC},k,p,\max}.
    \label{eq:erreichbarer_strombereich_stellkanal}
\end{equation}

Die strengen Ungleichungen verdeutlichen, dass der Sollwert nicht
unmittelbar auf einer Stellgrenze liegen darf. Sowohl in Richtung
kleinerer als auch in Richtung größerer Ströme ist eine Stellreserve
erforderlich. Diese Reserve wird zum Ausgleich thermischer
Widerstandsänderungen, elektrischer und mechanischer Toleranzen sowie
von Abweichungen der Speisespannung benötigt.

Für die erste Vorbemessung wird für Trafo~2 ein
Nennspannungsverhältnis von \SI{10}{\volt} zu
\SI{230}{\volt} angesetzt. Daraus folgt

\begin{equation}
    u_{\mathrm{T2}}
    =
    \frac{\SI{10}{\volt}}{\SI{230}{\volt}}
    \approx
    0{,}0435.
    \label{eq:vorbemessung_uebersetzung_trafo2}
\end{equation}

Bei einer vorläufig angesetzten primärseitigen Arbeitsspannung von

\begin{equation}
    U_{\mathrm{T2},1}
    =
    \SI{4.09}{\volt}
    \label{eq:vorbemessung_arbeitsspannung}
\end{equation}

ergibt sich für den idealisierten Transformator eine sekundärseitige
Spannung von

\begin{equation}
    U_{\mathrm{T2},2}
    =
    \frac{U_{\mathrm{T2},1}}{u_{\mathrm{T2}}}
    \approx
    \SI{94.1}{\volt}.
    \label{eq:vorbemessung_sekundaerspannung}
\end{equation}

Für die kleinen Abgänge kann der erforderliche sekundärseitige Strom
näherungsweise mit

\begin{equation}
    I_{\mathrm{T2},2}
    =
    u_{\mathrm{T2}}
    I_{\mathrm{T2},1}
    \label{eq:vorbemessung_sekundaerstrom}
\end{equation}

bestimmt werden. Der zugehörige Stellwiderstand im Arbeitspunkt folgt
unter Vernachlässigung der inneren Transformatorimpedanz aus

\begin{equation}
    R_{\mathrm{Stell},0}
    =
    \frac{U_{\mathrm{T2},2}}
         {I_{\mathrm{T2},2}}.
    \label{eq:vorbemessung_stellwiderstand}
\end{equation}

Tabelle~\ref{tab:vorbemessung_kleine_stellkanaele} zeigt die daraus
resultierenden idealisierten Arbeitspunkte.

\begin{table}[H]
    \centering
    \caption{Idealisierte Vorbemessung der Stellkanäle für
    \SI{40}{\ampere} und \SI{70}{\ampere}}
    \label{tab:vorbemessung_kleine_stellkanaele}
    \begin{tabular}{ccccc}
        \toprule
        \(I_{\mathrm{T2},1}\) &
        \(U_{\mathrm{T2},1}\) &
        \(U_{\mathrm{T2},2}\) &
        \(I_{\mathrm{T2},2}\) &
        \(R_{\mathrm{Stell},0}\) \\
        \midrule
        \SI{40}{\ampere} &
        \SI{4.09}{\volt} &
        \SI{94.1}{\volt} &
        \SI{1.74}{\ampere} &
        \SI{54.1}{\ohm} \\

        \SI{70}{\ampere} &
        \SI{4.09}{\volt} &
        \SI{94.1}{\volt} &
        \SI{3.04}{\ampere} &
        \SI{30.9}{\ohm} \\
        \bottomrule
    \end{tabular}
\end{table}

Die ideal übertragene Scheinleistung beträgt näherungsweise

\begin{equation}
    S_{\mathrm{T2}}
    \approx
    U_{\mathrm{T2},1} I_{\mathrm{T2},1}
    \approx
    U_{\mathrm{T2},2} I_{\mathrm{T2},2}.
    \label{eq:scheinleistung_trafo2}
\end{equation}

Damit ergeben sich ohne Auslegungsreserve etwa
\SI{164}{\volt\ampere} für einen \SI{40}{\ampere}-Kanal und
\SI{286}{\volt\ampere} für einen \SI{70}{\ampere}-Kanal. Diese Werte
stellen ausschließlich idealisierte Ausgangswerte dar. Für die
Komponentenauswahl müssen Transformatorverluste,
Kurzschlussimpedanz, Eigenerwärmung und die erforderliche
Stellreserve berücksichtigt werden.

Für die überschlägige Bewertung der Hochstromabgänge wird die
temperaturabhängige Widerstandsänderung des V2A-Grundlastpfads
herangezogen. Unter der Annahme einer konstanten Spannung gilt
näherungsweise

\begin{equation}
    \frac{I_{\mathrm{V2A}}(T)}
         {I_{\mathrm{V2A}}(T_0)}
    =
    \frac{1}
    {1+\alpha_{\mathrm{el}}(T-T_0)}.
    \label{eq:stromverhaeltnis_v2a_vorbemessung}
\end{equation}

Mit dem für die Voruntersuchung verwendeten Temperaturkoeffizienten

\begin{equation}
    \alpha_{\mathrm{el}}
    =
    \SI{0.001}{\per\kelvin},
\end{equation}

einer Ausgangstemperatur von \SI{20}{\degreeCelsius} und einer
angenommenen Temperatur von \SI{150}{\degreeCelsius} ergibt sich

\begin{equation}
    \frac{I_{\mathrm{V2A}}(\SI{150}{\degreeCelsius})}
         {I_{\mathrm{V2A}}(\SI{20}{\degreeCelsius})}
    =
    \frac{1}
    {1+\SI{0.001}{\per\kelvin}\cdot\SI{130}{\kelvin}}
    \approx
    0{,}885.
    \label{eq:stromverhaeltnis_v2a_150c}
\end{equation}

Der V2A-Teilstrom nimmt unter diesen vereinfachten Bedingungen somit
um etwa \SI{11.5}{\percent} ab. Würde der V2A-Widerstand im kalten
Zustand zunächst den vollständigen Nennstrom führen, ergäben sich die
in Tabelle~\ref{tab:orientierungswerte_nachfuehrstrom}
aufgeführten orientierenden Nachführströme.

\begin{table}[H]
    \centering
    \caption{Orientierungswerte für die thermisch erforderliche
    Stromnachführung der Hochstromabgänge}
    \label{tab:orientierungswerte_nachfuehrstrom}
    \begin{tabular}{ccc}
        \toprule
        Prüfstrom &
        relative Stromabnahme &
        orientierender Nachführstrom \\
        \midrule
        \SI{250}{\ampere} & \SI{11.5}{\percent} & \SI{28.8}{\ampere} \\
        \SI{400}{\ampere} & \SI{11.5}{\percent} & \SI{46.0}{\ampere} \\
        \SI{630}{\ampere} & \SI{11.5}{\percent} & \SI{72.5}{\ampere} \\
        \bottomrule
    \end{tabular}
\end{table}

Die Werte stellen keine abschließende Dimensionierung von Trafo~2
dar. Im entwickelten Konzept wird der V2A-Widerstand bereits im
kalten Zustand unterhalb des Sollstroms eingestellt. Trafo~2 führt
daher zusätzlich zum thermischen Nachführstrom einen anfänglichen
Grundstrom. Für die Auslegung ist folglich der größte während der
Prüfung auftretende Transformatorstrom maßgebend:

\begin{equation}
    I_{\mathrm{T2},1,k,p,\max}
    \geq
    I_{\mathrm{soll},k,p}
    -
    I_{\mathrm{V2A},k,p,\min}.
    \label{eq:auslegungsbedingung_trafo2_hochstrom}
\end{equation}

Gleichzeitig muss der minimale Strombeitrag des Transformatorzweigs
so klein sein, dass der Sollstrom im kalten Zustand nicht
überschritten wird:

\begin{equation}
    I_{\mathrm{T2},1,k,p,\min}
    \leq
    I_{\mathrm{soll},k,p}
    -
    I_{\mathrm{V2A},k,p,\mathrm{kalt}}.
    \label{eq:auslegungsbedingung_trafo2_minimalstrom}
\end{equation}

Der vorgesehene Arbeitspunkt des Stellwiderstands muss deshalb
zwischen seinen mechanischen und elektrischen Grenzen liegen:

\begin{equation}
    R_{\mathrm{Stell},\min}
    <
    R_{\mathrm{Stell},0}
    <
    R_{\mathrm{Stell},\max}.
    \label{eq:arbeitspunkt_stellwiderstand}
\end{equation}

Für Trafo~2 sind neben dem Übersetzungsverhältnis insbesondere die
Bemessungsleistung, die zulässigen Primär- und Sekundärströme, die
Kurzschlussimpedanz sowie die thermische Dauerbelastbarkeit
maßgebend. Für den Stellwiderstand müssen der minimale und maximale
Widerstandswert, der zulässige Dauerstrom, die Verlustleistung, die
mechanische Auflösung und die Verfahrgeschwindigkeit bekannt sein.

Da die Erwärmungsprüfung über mehrere Stunden durchgeführt wird,
müssen Trafo~2, Stellwiderstand, Schleifkontakt und sämtliche
elektrischen Verbindungen für Dauerbetrieb ausgelegt werden. Ihre
Eigenerwärmung darf den verfügbaren Stellbereich nicht unzulässig
verändern und die Temperaturmessung am eigentlichen Prüfling nicht
beeinflussen.

Die Vorbemessung zeigt, dass das Konzept mit technisch handhabbaren
sekundärseitigen Strömen und Widerständen grundsätzlich umsetzbar ist.
Eine abschließende Bewertung ist auf Grundlage der idealisierten
Berechnungen jedoch nicht möglich. Hierzu müssen die inneren
Transformatorimpedanzen, die temperaturabhängigen Widerstände, die
Quellenimpedanz und die Wechselwirkungen zwischen den Stellkanälen
berücksichtigt werden. Diese Zusammenhänge werden im folgenden
Kapitel als Simulationsmodell abgebildet und anhand definierter
Grenz- und Vergleichsfälle plausibilisiert.